\section*{Ethical Considerations}
This work aims to improve NLP capabilities for historically underserved languages in African communities. The downstream benchmarks we use are publicly available. The origin and licence of the monolingual pretraining text are not documented (Section~\ref{sec:data}); we therefore do not redistribute it. We do not introduce new data collection; potential biases in the source corpora may propagate to VEXMLM's representations.

\section{Limitations and Future Directions}
\label{sec:limitations}
Several limitations remain. First, the vocabulary expansion and tokenization strategy were developed specifically for Amharic and Tigrinya, and the approach has not been evaluated on other Ge'ez-script languages such as Tigre, Harari, or Classical Ge'ez, nor on other African languages. Second, the XLM-R baseline of the downstream task tables (Table~\ref{tab:downstream}, Table~\ref{tab:appendix_f1}) was run with a single seed, so its seed variance there is unknown; the VEXMLM results, and both models in the OOV and ablation tables, are averaged over five seeds. Third, the OOV word set used in Table~\ref{tab:oov} and Table~\ref{tab:ablation} is defined relative to the expanded tokenizer and includes words that XLM-R can represent but fragments more than VEXMLM, so it does not measure unseen words only. The ablation also does not separate the adaptation of the new embeddings from the additional training that continued pretraining provides, because no unexpanded model received the same continued pretraining. Fourth, the pretraining text is of undocumented origin, which limits the data statement we can give and prevents redistribution. Fifth, the proposed tokenizer remains primarily frequency-driven and does not explicitly model the rich morphological structure of Ge'ez-script languages. Finally, the QA development sets are small (\splitTigQAValidation{} TIGQA questions).

Our findings also suggest that effective vocabulary integration is as important as vocabulary expansion itself: the expanded vocabulary helps only after continued pretraining.

Future work will investigate vocabulary augmentation for additional Ge'ez-script languages, morphology-aware tokenization strategies, hybrid character-byte representations, a multi-seed XLM-R baseline, and the causes of the lower QA scores. Future studies should also examine alternative vocabulary integration methods and controlled comparisons with African-focused multilingual models such as AfroXLMR and AfriBERTa under a unified evaluation framework.
